\documentclass{beamer}
\usepackage{hyperref}
\usepackage[czech]{babel}

%\usetheme{Boadilla}
\usetheme{default}
\usecolortheme{seahorse}
\definecolor{myblue}{rgb}{0.2, 0.2, 0.6}
\setbeamertemplate{caption}{{\color{myblue}Obrázek:} \raggedright\insertcaption\par}
\setbeamertemplate{footline}[frame number] 

\newcommand{\authorname}{Dluhoš Matěj, Gaja Jan, Gajdošík Petr, Jurásek Petr, Phan Thanh Tú, Jan Rádl}
\newcommand{\authorsshort}{Dluhoš, Gaja, Gajdošík, Jurásek, Phan, Rádl}
\newcommand{\thesisname}{PDS: Python}

\title{\thesisname}
\author{\authorname}
\institute{Univerzita Palackého v Olomouci}
\date{\today}

\setbeamertemplate{navigation symbols}{}
\setbeamertemplate{headline}{}
\setbeamertemplate{footline}{
  \leavevmode%
  \hbox{%
  \begin{beamercolorbox}[wd=.4\paperwidth,ht=2.25ex,dp=1ex,center]{author in head/foot}%
    \usebeamerfont{author in head/foot}\authorsshort
  \end{beamercolorbox}%
  \begin{beamercolorbox}[wd=.3\paperwidth,ht=2.25ex,dp=1ex,center]{title in head/foot}%
    \usebeamerfont{title in head/foot}\thesisname
  \end{beamercolorbox}%
  \begin{beamercolorbox}[wd=.3\paperwidth,ht=2.25ex,dp=1ex,right]{date in head/foot}%
    \usebeamerfont{date in head/foot}\insertshortdate{}\hspace*{2em}
    \insertframenumber{} / \inserttotalframenumber\hspace*{2ex} 
  \end{beamercolorbox}}%
  \vskip0pt%
}

%\logo{\includegraphics[height=1cm]{UP_logo.png}}
\usepackage{tikz}
\addtobeamertemplate{headline}{}{%
    \begin{tikzpicture}[overlay, remember picture]
    	\ifnum\insertframenumber>1
        	\node[anchor=north east, inner sep=5pt] at (current page.north east) {\includegraphics[height=1cm]{obrazky/UP_logo.png}};
        \fi
    \end{tikzpicture}
}

\setbeamercolor{titlelike}{bg=,fg=}

\begin{document}

\begin{frame}
	\titlepage
\end{frame}

\begin{frame}{Obsah}
	\tableofcontents
\end{frame}

\section{Úvod do Pythonu s paralelismem}
\subsection{Historie a zaměření}
\begin{frame}{Historie a zaměření}
    \begin{columns}
        \column{0.55\textwidth}
        \begin{itemize}
            \item Vytvořen \textbf{Guidem van Rossumem} v roce \textbf{1991}.
            \item Jednoduchost a čitelnost kódu.
            \item Vysoká produktivita programátorů.
            \item Univerzální použití: od webových aplikací přes vědecké výpočty až po strojové učení.
            \item Python získal základní podporu paralelismu prostřednictvím vláken ve verzi 1.5.2 (1999).
        \end{itemize}
        
        \column{0.45\textwidth}
        \begin{figure}
        		\includegraphics[width=0.75\textwidth]{obrazky/Guido_van_Rossum.jpg}
        		\caption{Guido van Rossum}
      		\end{figure}
    \end{columns}
\end{frame}

\subsection{Rozdíl vláken a procesů}
\begin{frame}{Proces}
	\begin{itemize}
		\item Instance programu.
		\item \textbf{Klíčová fakta:}
			\begin{itemize}
				\item [\textendash] Využívá více jader a procesorů.
				\item [\textendash] Nový proces se spouští nezávisle na prvním procesu.
				\item [\textendash] Má oddělený paměťový prostor.
				\item [\textendash] Každý proces má svůj vlastní GIL.
				\item [\textendash] Ideální pro úlohy náročné na CPU.
				\item [\textendash] Procesy lze přerušit nebo ukončit.
			\end{itemize}
		\item \textbf{Nevýhody:}
			\begin{itemize}
				\item [\textendash] Spuštění procesu je pomalejší než spuštění vlákna.
				\item [\textendash] Vyšší paměťová náročnost.
				\item [\textendash] Složitější IPC.
			\end{itemize}
	\end{itemize}
\end{frame}

\begin{frame}{Vlákno}
	\begin{itemize}
		\item Entita v rámci procesu.
		\item \textbf{Klíčová fakta:}
			\begin{itemize}
				\item [\textendash] V rámci jednoho procesu může být spuštěno více vláken.
				\item [\textendash] Paměť je sdílena mezi všemi vlákny.
				\item [\textendash] Spuštění vlákna je rychlejší než spuštění procesu.
				\item [\textendash] Skvělé pro úlohy závislé na I/O.
				\item [\textendash] Lehká paměťová náročnost.
			\end{itemize}
		\item \textbf{Nevýhody:}
			\begin{itemize}
				\item [\textendash] Jedno GIL pro všechna vlákna.
				\item [\textendash] Multithreading nemá efekt u úloh náročných na CPU.
				\item [\textendash] Vlákna nelze přerušit a ukončit.
				\item [\textendash] Zvýšený potenciál vzniku chyby souběhu.
			\end{itemize}
	\end{itemize}
\end{frame}

\subsection{GIL - Global Interpreter Lock}
\begin{frame}{Co to je GIL a proč je potřeba?}
	\textbf{Co to je GIL?}
	\begin{itemize}
		\item GIL je mutex umožňující pouze jednomu vláknu mít kontrolu nad Python interpreterem.
	\end{itemize}
	\textbf{Proč je potřeba?}
	\begin{itemize}
		\item Správa paměti v CPythonu není bezpečná pro více vláken.
		\item Python využívá počítání referencí.
		\item Problém nastane, když dvě vlákna současně zvyšují nebo snižují referenci.
		\item To vede k úniku nebo k nesprávnému uvolnění paměti.
	\end{itemize}
\end{frame}

\begin{frame}{Potřebujeme vůbec GIL?}
	\textbf{Je tedy potřeba?}
	\begin{itemize}
		\item V brzkých dnech Pythonu usnadnil integraci nativního kódu v C.
		\item Vyřešil problémy s přístupem ke sdíleným datům.
		\item Optimalizoval výkon na systémech s jedním jádrem.
	\end{itemize}
	\textbf{Ale...}
	\begin{itemize}
		\item Omezuje paralelní výkon na vícejádrových procesorech.
		\item Vede k neefektivnímu využití vláken u CPU-bound úloh.
	\end{itemize}
	\textbf{\hypersetup{urlcolor=blue} \href{https://peps.python.org/pep-0703/}{PEP 703}} - Making the Global Interpreter Lock Optional in CPython.
\end{frame}


%\subsection{Paralelismus v Pythonu}
%\begin{frame}{Paralelismus v Pythonu}
%	\begin{itemize}
%		\item Python získal základní podporu paralelismu prostřednictvím vláken ve verzi 1.5.2 (1999).
%		\item Zaveden \textbf{GIL} (Global Interpreter Lock)
%		\begin{itemize}
%			\item [\textendash] Usnadnilo integraci nativního kódu v C.
%			\item [\textendash] Vyřešilo problémy s přístupem ke sdíleným datům ve více vláknech.
%			\item [\textendash] Optimalizovalo výkon na systémech s jedním jádrem.
%			\item [\textendash] Omezuje paralelní výkon na vícejádrových procesorech.
%			\item [\textendash] Vede k neefektivnímu využití vláken u \textbf{CPU-bound} úloh.
%			\item [\textendash] \hypersetup{urlcolor=blue} \href{https://peps.python.org/pep-0703/}{PEP 703} - Making the Global Interpreter Lock Optional in CPython 
%		\end{itemize}
%	\end{itemize}
%\end{frame}

%\subsection{Proč odstranit GIL?}
%\begin{frame}{Proč odstranit GIL?}
%	\begin{itemize}
%		\item Většina zařízení dnes obsahuje vícejádrové procesory, kde GIL představuje významné omezení.
%		\item Efektivnější paralelní výpočty, například v oblastech umělé inteligence nebo numerických simulacích.
%		\item Výzvy při odstranění GIL:
%		\begin{itemize}
%			\item [\textendash] Potřeba bezpečného a efektivního spravování sdílené paměti (například více zamykání nebo jiných synchronizačních mechanismů).
%			\item [\textendash] Zvýšení složitosti implementace interpretu Pythonu.
%		\end{itemize}
%	\end{itemize}
%\end{frame}


\section{Synchronizační nástroje standardní knihovny}
\subsection{Knihovna threading}
\begin{frame}{Knihovna threading}


\end{frame}

\subsection{Knihovna multiprocessing}
\begin{frame}{Knihovna multiprocessing}

\end{frame}

\subsection{Knihovna concurrent.futures}
\begin{frame}{Knihovna concurrent.futures}

\end{frame}

\subsection{Knihovna asyncio}
\begin{frame}{Knihovna asyncio}

\end{frame}

\section{Nástroje externích knihoven}
\begin{frame}{Nástroje externích knihoven}

\end{frame}

\section{Nástroje pro distribuované systémy}
\subsection{Knihovna socket}
\begin{frame}{Knihovna socket}

\end{frame}

\subsection{Knihovna http.server}
\begin{frame}{Knihovna http.server}

\end{frame}

\subsection{RabbitMQ}
\begin{frame}{RabbitMQ}

\end{frame}

\subsection{Další nástroje}
\begin{frame}{Další nástroje}

\end{frame}

\end{document}